% =====================================================================================
% Chapter 9 · Geo lift — synthetic control
%
% Built from notebooks/07_geo_lift_synthetic_control.ipynb (the body) and
% notebooks/07b_geo_lift_real_data.ipynb (§9.10). NOT ONE NUMBER IN THIS FILE IS TYPED
% BY HAND: every figure in the prose is a macro from build/macros.tex, generated by
% cmp.report from the executed notebooks. The only literals below are (i) the constants
% of the data-generating model in §9.2, which are *definitions* rather than results, and
% (ii) the decision-rule thresholds (0.9 / 0.5) and Abadie's 5x pre-fit filter, which are
% *conventions* stated in the text, not measurements.
% =====================================================================================
\chapter{Geo lift: synthetic control}
\label{ch:geo}

\begin{quote}\small
\emph{One market got the campaign. Thirty did not. There was no coin flip, and the
counterfactual — what that one market would have sold with no campaign — was never
observed and never will be. This chapter builds it, prices it, and then discovers that the
Bayesian interval around it is badly too narrow. The diagnosis is precise, and it is the
spine of this book: the likelihood was wrong, not the paradigm.}
\end{quote}

% =====================================================================================
\section{The decision}
\label{sec:sc:decision}

A television campaign ran in a single metropolitan area --- a \emph{designated market area},
the standard unit of media buying --- and nowhere else. Sales in that metro have risen since
launch. The proposal on the table is a national rollout, and the campaign that produced the
evidence cost \keur{\nbSevenCost}. The question is not whether sales rose. They did. The
question is whether the \emph{campaign} raised them, and by enough to pay for itself.

A wrong answer is expensive in both directions, and asymmetrically so. Credit the campaign
with a lift the market would have delivered anyway and the firm rolls a dud out nationally at
many multiples of the test budget. Miss a real lift and it walks away from a working channel
that a competitor will find. What makes the problem hard is stated in one sentence: \textbf{there
is no control group}. Nobody flipped a coin. The treated metro was chosen, and the rest of the
country was left dark by default.

The tempting shortcuts follow immediately, and all of them fail. A before-and-after comparison
within the treated metro books the shared trend, the season and the state of the economy as
campaign effect. A comparison against the average of the other markets books the level gap
between a large metro and an average one. Both failures are quantified in \cref{tab:sc:naive},
and both have the same root: every market in the country is riding the same waves, and the
estimator has to know which part of the treated market's rise was the wave.

What is needed is the one quantity that is never observed: the treated market's sales
\emph{had the campaign not run}. \citet{abadie2003} proposed constructing it as a weighted blend
of the untreated markets --- the \emph{donor pool} --- with the weights chosen so that the blend
tracks the treated market before the campaign. If a mix of, say, sixty per cent of one metro and
forty per cent of two others matched the treated market week for week for forty weeks, then, so
long as those markets keep moving together, that same mix is the best available guess of what the
treated market would have done afterwards. The post-launch gap between the market and its
synthetic double is the estimated lift. This chapter builds that estimator classically, rebuilds
it as a posterior, and puts the two on the same estimand.

% =====================================================================================
\section{The data-generating model}
\label{sec:sc:dgm}

A causal method cannot be graded on real data, because on real data the counterfactual is
missing --- that is the entire problem. So the method is first graded against a world whose
answer is known by construction, and only then trusted on a world whose answer is not.
\Cref{sec:sc:real} performs the second half of that bargain on a real, randomized geo
experiment; this section builds the first.

The simulator produces a weekly sales panel: \nbSevenNDmas{} markets, indexed
$i = 0, \dots, \nbSevenNDmas-1$, over \nbSevenNWeeks{} weeks $t = 0, \dots, \nbSevenNWeeks-1$.
Three latent factors are \emph{shared by every market} --- a secular trend $g_t$, a seasonal
cycle $s_t$, and a macroeconomic random walk $m_t$:
%
\begin{align}
  g_t &= 0.4\,t, \label{eq:sc:trend}\\
  s_t &= 8\sin\!\Big(\tfrac{2\pi t}{26}\Big) + 4\sin\!\Big(\tfrac{2\pi t}{13}\Big),
        \label{eq:sc:season}\\
  m_t &= \sum_{u \le t} \eta_u, \qquad \eta_u \sim \mathcal N\!\left(0,\ 1.2^{2}\right).
        \label{eq:sc:macro}
\end{align}
%
Each market carries its own baseline $c_i \sim U(80, 140)$ and its own sensitivity to each
factor --- its \emph{loadings} $\lambda_i = (\lambda_{i1}, \lambda_{i2}, \lambda_{i3})$, each drawn
$U(0.6, 1.4)$ --- and observed sales, in thousands of euros, are
%
\begin{equation}
  Y_{it} \;=\; c_i \;+\; \lambda_{i1}\,g_t \;+\; \lambda_{i2}\,s_t \;+\; \lambda_{i3}\,m_t
           \;+\; \varepsilon_{it},
  \qquad \varepsilon_{it} \sim \mathcal N\!\left(0,\ 3^{2}\right).
  \label{eq:sc:dgp}
\end{equation}
%
The campaign is planted in market $0$ alone, from week \nbSevenLaunchWeek{} onwards, as a
proportional lift of \pc{\nbSevenLiftPct} on that market's untreated level:
%
\begin{equation}
  \Delta_t \;=\; 0.12 \,\big(c_0 + \lambda_{01}\,g_t + \lambda_{02}\,s_t + \lambda_{03}\,m_t\big)
  \cdot \mathbf 1[\,t \ge \nbSevenLaunchWeek\,],
  \qquad Y_{0t} \;\mathrel{+}=\; \Delta_t .
  \label{eq:sc:lift}
\end{equation}
%
Summed over the \nbSevenNPost{} post-launch weeks, the planted total lift is
\keur{\nbSevenTrueTotal}, an average of \keur{\nbSevenTrueAvg} per week. That number is the
grade every estimator in this chapter is marked against, and it is the last time in the chapter
that anything is known for certain.

\Cref{fig:sc:panel} shows the panel. The macro random walk $m_t$ of \cref{eq:sc:macro} is both
the villain and the hero of the story. It is the villain because it moves every market together,
so any before-and-after comparison will book its drift as campaign effect. It is the hero because
it is \emph{shared}: since all markets load on the same $(g_t, s_t, m_t)$, a convex combination
of donor markets whose blended loadings match $(\lambda_{01}, \lambda_{02}, \lambda_{03})$
reproduces market $0$'s untreated path exactly, up to idiosyncratic noise. That is what makes the
effect identifiable at all. It is also --- and this only becomes visible in \cref{sec:sc:compare}
--- the reason the Bayesian interval will fail, because $m_t$ is a random walk and what the
weights fail to match of it does not average away.

\input{build/floats/nb07_panel}

% =====================================================================================
\section{Identification}
\label{sec:sc:ident}

\subsection*{The estimand}

In the potential-outcome notation of \citet{rubin1974} and \citet{imbens2015}, write
$Y_{0,t}(1)$ for market $0$'s week-$t$ sales when the campaign runs and $Y_{0,t}(0)$ for its
sales when it does not. Only the first is observed. The causal effect in week $t$ is their
contrast, and the object the business is buying is its sum over the post-launch window:
%
\begin{equation}
  \tau_t \;=\; Y_{0,t}(1) - Y_{0,t}(0), \qquad t \ge T_0,
  \qquad\qquad
  \tau \;=\; \sum_{t \ge T_0} \tau_t ,
  \label{eq:sc:estimand}
\end{equation}
%
with $T_0 = \nbSevenLaunchWeek$ the launch week. Everything in this chapter --- the classical
estimator, the posterior, both intervals, and the euro decision --- targets the single scalar
$\tau$ of \cref{eq:sc:estimand}. Keeping the estimand fixed is what makes the comparison in
\cref{sec:sc:compare} a comparison rather than a change of subject.

\subsection*{The estimator}

Synthetic control fills the hole in \cref{eq:sc:estimand} with a weighted average of the
$J = \nbSevenNDonors$ donor markets:
%
\begin{equation}
  \widehat{Y_{0,t}(0)} \;=\; \sum_{j=1}^{J} w_j\, Y_{jt},
  \qquad
  \hat\tau_t \;=\; Y_{0t} - \sum_{j=1}^{J} w_j\, Y_{jt},
  \label{eq:sc:gap}
\end{equation}
%
with the weights confined to the unit simplex,
$\simplex = \{\,w : w_j \ge 0,\ \sum_j w_j = 1\,\}$. The simplex is not decoration. Left
unconstrained --- ordinary regression on the donors --- the ``synthetic'' market may sit outside
the range of every real market, which is extrapolation wearing the costume of a fit: it matches
the pre-period beautifully and then diverges after it. Restricting $w$ to $\simplex$ keeps the
synthetic inside the convex hull of the donor pool, and that restriction is what makes the
post-launch projection a defensible one \citep{abadie2010, abadie2021}. \Cref{sec:sc:real} shows
a real market that lies outside its own donor hull and what the method does about it.

\subsection*{The assumptions}

Identification of \cref{eq:sc:estimand} by the estimator \cref{eq:sc:gap} rests on five
assumptions. They are stated here, and each is either tested or explicitly declared untestable
later in the chapter --- an assumption with no accompanying test is a wish.

\begin{assumption}[Convex hull]\label{as:sc:hull}
The treated market lies inside the convex hull of the donor pool: there exists
$w \in \simplex$ whose blended factor loadings equal
$(\lambda_{01}, \lambda_{02}, \lambda_{03})$. The simplex constraint enforces this by refusing
to extrapolate; the diagnostic is a pre-period fit that the weights cannot achieve.
\end{assumption}

\begin{assumption}[Good pre-fit]\label{as:sc:prefit}
The synthetic tracks the treated market \emph{before} launch. Formally, the pre-period RMSE of
\cref{eq:sc:gap} is small relative to the effect the design must detect. This is the
precondition on which every subsequent number depends: a synthetic that cannot reproduce the
treated market's past cannot be trusted with its counterfactual future. It is tested as a
pass/fail gate in \cref{sec:sc:valid}.
\end{assumption}

\begin{assumption}[No anticipation]\label{as:sc:anticip}
Nothing shifts sales in the treated market \emph{before} $T_0$ --- no leaked creative, no
pre-buying, no mis-dated launch. Tested by placebo-in-time in \cref{sec:sc:wrong}.
\end{assumption}

\begin{assumption}[Stable factor structure]\label{as:sc:stable}
The co-movement that held in the pre-period continues to hold after it: the loadings the weights
matched on weeks $0$ to $T_0 - 1$ still govern weeks $T_0$ onward. A long pre-period makes the
match harder to achieve by accident; placebo-in-time probes it.
\end{assumption}

\begin{assumption}[No interference]\label{as:sc:sutva}
The campaign in market $0$ does not change sales in any donor market --- SUTVA
\citep{imbens2015}, applied to geographies. This is the assumption that no placebo can test,
because spillover contaminates every placebo refit in exactly the same way. It is stress-tested
by simulation in \cref{sec:sc:wrong}, where the news is better than it sounds: spillover
\emph{attenuates}.
\end{assumption}

\begin{remark}
A sixth condition --- that nothing else hit market $0$ at exactly $T_0$ --- is genuinely
untestable from the panel and can be discharged only by institutional knowledge. Say so, rather
than dress a diagnostic in its clothes. \Cref{sec:sc:real} is instructive here: on a randomized
geo experiment, that assumption is guaranteed in expectation by the coin flip, and the difference
between the two situations is exactly the value of an experiment.
\end{remark}

% =====================================================================================
\section{The classical baseline}
\label{sec:sc:classical}

Before any sampler, the discipline of this book is to ask what a competent analyst produces
\emph{without} one. In this chapter that question has an unusually clean answer, because the
classical estimator \emph{is} the method: \citeauthor{abadie2003}'s synthetic control was born
frequentist, and the Bayesian model of \cref{sec:sc:bayes} is its re-housing, not its rival.

\subsection*{The estimator, without a likelihood}

Choose the weights that best reproduce the treated market's pre-launch path, subject to
\cref{as:sc:hull}'s simplex:
%
\begin{equation}
  \hat w \;=\; \argmin_{w \,\in\, \simplex}\ \sum_{t < T_0}
    \Big( Y_{0t} - \sum_{j=1}^{J} w_j\, Y_{jt} \Big)^{2}.
  \label{eq:sc:slsqp}
\end{equation}
%
This is constrained least squares --- a quadratic objective on a convex set --- solved
numerically by sequential least-squares quadratic programming. It is deterministic, it runs in
milliseconds, and it estimates nothing else: no variance parameter, no prior, no chain. The
counterfactual and the effect follow by substitution into \cref{eq:sc:gap}. Because it is so
cheap, this same fitter carries every robustness check in the chapter --- the leave-one-out
refits, the placebo distributions, the launch-date sweep, the spillover stress test.

Fitted to the pre-launch weeks, it puts material weight on \nbSevenClNMaterial{} of the
\nbSevenNDonors{} donors --- an effective donor count of \nbSevenClEffDonors{}, measured by the
inverse Herfindahl index $1/\sum_j w_j^2$ --- and tracks the treated market to a pre-launch RMSE
of \keur{\nbSevenClPreRmse} (\cref{fig:sc:sc}). \Cref{as:sc:prefit} is satisfied, so the
post-launch gap can be read. It comes to \keur{\nbSevenClAvg} per week and \keur{\nbSevenClTotal}
over the window, against a planted truth of \keur{\nbSevenTrueTotal} --- a relative error of
$\nbSevenClErrPct\,\%$.

\input{build/floats/nb07_sc}

\subsection*{Where an interval comes from when there is no likelihood}

Nowhere, by the usual route. There is no likelihood, so there is no standard error to read off a
fit; and a regression standard error, even if one were bolted on, would answer the wrong
question, because standard errors describe variability across many independent treated units and
there is exactly one. \citeauthor{abadie2010}'s answer --- the inferential heart of the method
--- is \emph{permutation inference} \citep{abadie2010, abadie2015}, in the tradition of
\citet{fisher1935}:

\begin{quote}
Under the null that the campaign did nothing, the treated market is just another market. So
refit the identical estimator on each donor \emph{as if it had been treated}, collect the gaps
those placebo fits produce --- that is the distribution of gaps the world hands you from noise
alone --- and ask where the real one ranks.
\end{quote}

Two refinements matter, and both are \citeauthor{abadie2010}'s. First, \emph{discard placebos
with a poor pre-fit}: a donor whose own synthetic cannot track it before launch produces a
meaningless gap that would only pollute the reference distribution. The conventional cutoff, used
here, is a pre-period RMSE more than five times the treated unit's; \nbSevenPlaceboKept{} of the
\nbSevenNDonors{} donors survive it. Second, \emph{rank on the RMSE ratio rather than the raw
gap}: post-period RMSE divided by pre-period RMSE is scale-free, so it is fair across large and
small markets, and it refuses to reward a placebo that simply fits badly everywhere, since such a
unit has a high numerator \emph{and} a high denominator.

On that statistic the treated market ranks \nbSevenTreatedRank{} of \nbSevenNUnits{}, with a
ratio of \nbSevenTreatedRatio{} against a placebo median of \nbSevenPlaceboMedianRatio{}, for
$p = \nbSevenPRatio$ (\cref{fig:sc:placebo}). The raw-gap permutation gives
$p = \nbSevenPGap$; the two differ by a few hundredths because of bookkeeping --- the raw-gap
version is an add-one permutation that excludes the treated unit from its own reference set,
while the ratio statistic is a rank that includes it --- and they agree on everything that
matters. The verdict is also stable in the one arbitrary knob: tightening the pre-fit filter to
$2\times$ gives $p = \nbSevenPFilterTwo$ and loosening it to $20\times$ gives
$p = \nbSevenPFilterTwenty$, the latter simply because a loose filter admits broken synthetics
and dilutes the test.

\input{build/floats/nb07_placebo}

\subsection*{An interval by inversion}

The same null yields an interval. The placebo totals map out how far this estimator's
\nbSevenNPost-week total wanders when the true effect is zero, so subtracting that distribution's
upper and lower percentiles from the estimate inverts the test into a 90\,\% \emph{randomization
interval}:
%
\begin{equation}
  \Big[\ \hat\tau - q_{0.95}\big(\tau^{\text{placebo}}\big),\ \
         \hat\tau - q_{0.05}\big(\tau^{\text{placebo}}\big)\ \Big]
  \;=\; \big[\,\keur{\nbSevenRiLo},\ \keur{\nbSevenRiHi}\,\big],
  \label{eq:sc:ri}
\end{equation}
%
valid under a constant-effect null: the placebos supply the \emph{spread} of the estimator's
error, not its location. The planted \keur{\nbSevenTrueTotal} falls inside it. This interval
never assumed a shape for the errors --- it went and measured one, from markets known to be
untreated. Hold on to that sentence; \cref{sec:sc:compare} is going to cash it.

What this classical apparatus cannot produce, at any level of effort, is
$\Prob(\tau > \text{cost})$. A $p$-value ranks data against a null and attaches no probability to
a hypothesis about the effect; \cref{eq:sc:ri} is a statement about a procedure, not about $\tau$.
The euro rule of \cref{sec:sc:euro} is nothing but that probability. That is the gap the next
section exists to fill.

% =====================================================================================
\section{The Bayesian model}
\label{sec:sc:bayes}

The natural probability distribution over the simplex is the Dirichlet: its draws are
non-negative and sum to one by construction, so it says ``the weights live on $\simplex$'' in the
language of priors rather than constraints. Fitting the pre-period gives a posterior over the
weights, hence over the whole counterfactual path, hence over $\tau$ --- \cref{eq:sc:slsqp}'s
objective with uncertainty attached.

The model, fitted on the pre-period only ($t < T_0$), is
%
\begin{align}
  Y_{0t} \mid w, \sigma &\;\sim\; \mathcal N\Big(\textstyle\sum_{j=1}^{J} w_j\, Y_{jt},\ \
                                   \sigma^{2}\Big), \qquad t < T_0, \label{eq:sc:lik}\\
  w &\;\sim\; \text{Dirichlet}(\mathbf 1_J), \label{eq:sc:prior_w}\\
  \sigma &\;\sim\; \text{HalfNormal}(5). \label{eq:sc:prior_s}
\end{align}
%
$\text{Dirichlet}(\mathbf 1_J)$ is flat over the simplex: no donor is favoured a priori. The
prior on $\sigma$ is weakly informative and tied to the units of the data --- sales are in
thousands of euros, weekly levels sit in the €80--170\,k band, and the idiosyncratic noise of
\cref{eq:sc:dgp} has standard deviation €3\,k, so a scale of €5\,k is wide enough to accommodate
any plausible blend error and tight enough to exclude an absurd one. As throughout this book, the
prior is doing very little work here; \cref{eq:sc:lik} --- the likelihood --- is doing all of it,
and \cref{sec:sc:compare} is about precisely that.

The counterfactual for a post-launch week is the posterior \emph{predictive} blend,
%
\begin{equation}
  \widehat{Y_{0t}(0)} \;=\; \sum_j w_j\, Y_{jt} \;+\; \varepsilon_t,
  \qquad \varepsilon_t \sim \mathcal N(0, \sigma^2),
  \label{eq:sc:ppc}
\end{equation}
%
and the $\varepsilon_t$ term earns its place: dropping it --- reporting the band of the mean
blend alone --- would carry weight uncertainty but not the treated market's own noise, and would
understate the per-week band. Posterior draws of $\tau$ follow by summing
$Y_{0t} - \widehat{Y_{0t}(0)}$ over the window.

Inference is by the No-U-Turn sampler \citep{hoffman2014} as implemented in PyMC
\citep{salvatier2016}: \nbSevenNDraws{} post-warmup draws, with $\hat R = \nbSevenRhat$, a
minimum bulk effective sample size of \nbSevenEss{} \citep{vehtari2021}, and \nbSevenDivergences{}
divergent transitions. The chains have converged; nothing that follows is a sampler artefact.

The fit puts weight on many more donors than \cref{eq:sc:slsqp} did --- \nbSevenBayesEffDonors{}
effective donors against the classical \nbSevenClEffDonors{} (\cref{tab:sc:weights}) --- because
a flat Dirichlet regularizes the blend. That is a real trade and worth naming: the regularized
blend is more stable out of sample, and less interpretable as a business object. ``Our metro
behaves like this specific mix of four named metros'' is a sentence a media planner can act on;
``like a diffuse average of fifteen'' is not.

\input{build/tables/nb07_weights}

The posterior mean total lift is \keur{\nbSevenBayesTotal}, with a 90\,\% credible interval of
$[\keur{\nbSevenBayesLo},\ \keur{\nbSevenBayesHi}]$ --- width \keur{\nbSevenBayesWidth}, against
the classical interval's \keur{\nbSevenRiWidth}. The planted \keur{\nbSevenTrueTotal} sits just
\emph{outside} that band. One panel is one draw of the world and a single miss proves nothing, so
that observation is not yet a finding. \Cref{sec:sc:valid} turns it into one.

% =====================================================================================
\section{Diagnostics and validation}
\label{sec:sc:valid}

\subsection*{The pre-fit gate}

\Cref{as:sc:prefit} is a gate, not a formality, and it is made concrete rather than eyeballed:
the pre-period RMSE must fall below a third of the weekly lift the design is meant to detect. (The
benchmark is the \emph{estimated} weekly lift, not the planted one --- a diagnostic that consults
the truth is not available on real data and therefore not available here.) The fit gives a
pre-RMSE of \keur{\nbSevenBayesPreRmse} against a bar of \keur{\nbSevenGateRef}:
\textbf{\nbSevenGate}. The pre-launch gap averages \keur{\nbSevenPreGapMean}, which is the shape
a good synthetic makes --- hugging zero before launch, departing after it (\cref{fig:sc:gap}).

\input{build/floats/nb07_gap}

\subsection*{Donor robustness}

If the estimate hangs on one market, it is not an estimate but an anecdote. Dropping each of the
top-weighted donors in turn and refitting moves the total only within
\keur{\nbSevenLooLo}--\keur{\nbSevenLooHi}, against the same fitter's all-donor baseline of
\keur{\nbSevenClTotal}. No single donor drives the result.

\subsection*{Calibration across fresh panels}

This is the check that pays for the simulator, and the one that turns the chapter. A single panel
is a single draw of the world; the honest question about an interval is not whether it happens to
cover this time but how often it covers. So the whole procedure --- simulate, fit, interval --- is
re-run on \nbSevenNPanels{} fresh panels drawn from \cref{eq:sc:dgp}, and the coverage of each
nominal 90\,\% interval is counted.

Three results, and they do not point the same way.

\begin{enumerate}[leftmargin=*]
  \item \textbf{Recovery is acceptable, precision is not.} The mean estimation error across panels
        is $\nbSevenSeedBias$ (EUR 000), or $\nbSevenSeedBiasPct\,\%$ of the truth --- a modest
        bias --- but the seed-to-seed standard deviation of that error is \keur{\nbSevenSeedSd}.
        Synthetic control here is a high-variance estimator, not a conservative one.
  \item \textbf{The per-week bands are honest.} The 90\,\% posterior-predictive band covers the
        true weekly path \pc{\nbSevenWeekCov} of the time --- essentially nominal. The
        observation-noise term of \cref{eq:sc:ppc} earns this.
  \item \textbf{The total-lift interval is not.} The 90\,\% credible interval on $\tau$ covers the
        planted total on only \pc{\nbSevenTotalCov} of panels. A nominal ninety that delivers
        \pc{\nbSevenTotalCov} is not a rounding error; it is a broken instrument.
\end{enumerate}

The classical arm, run on the very same panels, covers at \pc{\nbSevenRiCov} --- it keeps its
promise. \Cref{tab:sc:compare} records both. The rest of this chapter is about why.

% =====================================================================================
\section{Point estimate versus posterior}
\label{sec:sc:compare}

\input{build/tables/nb07_compare}

\subsection*{On the number, they agree}

\Cref{tab:sc:compare} puts both arms on the estimand of \cref{eq:sc:estimand}. The point estimates
differ by \keur{\nbSevenEstDiff} --- \pc{\nbSevenEstDiffPct} of the effect --- and both sit below
the planted truth. If anything the classical point estimate is the closer of the two, because the
Dirichlet prior spreads weight across donors and a regularized blend is a slightly more
conservative one. The Bayesian layer did not buy a better point estimate, and it could not have:
the causal work was done by the identification of \cref{sec:sc:ident} and by the simplex, and both
arms share it. This is the ordinary case in this book, and it has a name --- the Bernstein--von
Mises phenomenon \citep{gelman2013}: with a well-specified likelihood and enough data, the
posterior concentrates where the maximum-likelihood estimate does, and the two paradigms return
the same number.

\subsection*{On the uncertainty, they do not --- and the classical arm wins}

Now the width. Three quantities claim to price the same object, the standard deviation of the
\nbSevenNPost-week total, and they are set against each other in \cref{tab:sc:scales}. The
posterior says \keur{\nbSevenBayesSd}. The placebo distribution says \keur{\nbSevenPlaceboSd}.
And the referee --- the \emph{actual} spread of the estimation error across the
\nbSevenNPanels{} fresh panels, which only a simulator can supply --- says \keur{\nbSevenSeedSd}.

\input{build/tables/nb07_scales}

The posterior underprices the truth by a factor of \nbSevenUnderprice{}. The design-based placebo
scale lands within \pc{\nbSevenPlaceboSdErrPct} of it. That is the whole finding, and it is
measured, not argued: the coverage column of \cref{tab:sc:compare} is its consequence.

\subsection*{Why: a misspecified likelihood, not a failed paradigm}

The cause is a single modelling choice, and stating it precisely is the most important paragraph
in this chapter.

The likelihood \cref{eq:sc:lik} treats the weekly errors as independent. Under independence, the
uncertainty of an $n$-week total scales as
%
\begin{equation}
  \sd\Big(\sum_{t \in \text{post}} e_t\Big) \;=\; \sqrt{n}\,\sigma
  \qquad \text{(iid errors)} .
  \label{eq:sc:sqrtn}
\end{equation}
%
But the blend's error is not iid. Whatever share of the treated market's factor loadings the
weights fail to match leaks the \emph{shared} factors straight into the gap:
%
\begin{equation}
  e_t \;\approx\; \big(\lambda_0 - \textstyle\sum_j w_j \lambda_j\big)^{\!\top}
        (g_t,\, s_t,\, m_t) \;+\; \varepsilon_{0t},
  \label{eq:sc:leak}
\end{equation}
%
and $m_t$ is a random walk (\cref{eq:sc:macro}). The mismatch term is therefore a persistent,
slowly wandering offset --- not fresh noise every week. For correlated errors the variance of a
sum collects every covariance term,
%
\begin{equation}
  \Var\Big(\sum_{t=1}^{n} e_t\Big)
    \;=\; \sigma^{2}\Big( n \;+\; 2\sum_{k=1}^{n-1}(n-k)\,\rho_k \Big)
    \;\xrightarrow[\ \rho_k \to 1\ ]{}\; n^{2}\sigma^{2} \;\gg\; n\,\sigma^{2},
  \label{eq:sc:n2}
\end{equation}
%
where $\rho_k$ is the lag-$k$ autocorrelation of the error. \Cref{eq:sc:n2} is the entire
mechanism. \emph{Per-week bands need only the marginal variance} $\sigma^2$, which
\cref{eq:sc:ppc} gets roughly right --- which is exactly why they cover at \pc{\nbSevenWeekCov}.
\emph{The total needs the covariances too}, and \cref{eq:sc:lik} sets every $\rho_k$ to zero.

\Cref{fig:sc:acf} measures $\rho_k$ rather than assuming it, on placebo worlds where the true
lift is zero and the post-launch gap therefore \emph{is} the error. It is
$\rho_1 = \nbSevenRhoOne$ and still $\nbSevenRhoEight$ at lag eight: the error is almost perfectly
persistent. Substituting into \cref{eq:sc:n2} reproduces the observed underpricing without any
further appeal to the data.

\input{build/floats/nb07_acf}

One diagnostic deserves emphasis because it is the one that would have misled a careful analyst.
The \emph{in-sample} pre-period residual of the fitted model shows only modest, sign-flipping
autocorrelation --- $\nbSevenAcfPreOne$, $\nbSevenAcfPreTwo$ and $\nbSevenAcfPreThree$ at lags one
to three, against a noise band of roughly $\pm\nbSevenAcfNoiseBand$. The fit \emph{whitens} what
it can see, because the weights were chosen to absorb precisely the in-sample structure. Eyeballing
the pre-fit residual would therefore have raised no alarm at all. The persistence is a property of
the \emph{out-of-sample} error, and only a design-based procedure --- one that refits the estimator
on units known to be untreated --- ever gets to look at it.

So: the Bayesian arm lost here, and it lost for a reason that is fully diagnosed and entirely
repairable. It is \textbf{not} a failure of Bayesian inference. It is a likelihood asserting
independence over a quantity whose errors are serially correlated, and then being believed. Two
principled repairs exist, and neither is ``stop using Bayes''. The first is to give the model the
error structure it actually has --- an AR(1) or random-walk residual --- so that it can see the
persistence itself; this is the route taken by Bayesian structural time-series approaches to the
same problem \citep{brodersen2015}. The second is to correct the posterior for misspecification
directly: \citet{mueller2013} shows that a Bayesian posterior computed under a misspecified
likelihood can be rescaled by a sandwich covariance matrix so that its intervals recover their
frequentist coverage --- exactly the repair this chapter's situation calls for, since the
posterior's \emph{centre} is fine and only its \emph{scale} is wrong.

The classical robustness toolkit exists precisely to be agnostic about the error structure.
Clustered and HAC standard errors \citep{newey1987, bertrand2004}, the bootstrap, and
randomization inference all decline to assume a shape for the errors and go and measure one
instead. \Citet{bertrand2004} made the same point about difference-in-differences, and it cost the
literature a decade of overstated significance to absorb: serially correlated errors, aggregated
over time, destroy naive standard errors. Here the same lesson arrives in Bayesian dress. That is
why the design-based interval wins in this chapter, and it is the honest reason --- not because
frequentism is deeper, but because it committed to less.

\subsection*{What the posterior did buy}

One thing, and it is decisive. $\Prob(\tau > \text{cost})$ exists in exactly one of the two
columns of \cref{tab:sc:compare}. The classical entry is not ``small'' or ``hard to compute'';
it is \emph{undefined}, and that is a literal statement rather than a rhetorical concession. A
$p$-value ranks data against a null. No amount of placebo refitting turns a rank into a
probability that the campaign paid. \Cref{sec:sc:euro}'s rule is nothing but that probability, and
so the decision quantity lives only in the Bayesian arm.

Which sets up the uncomfortable part, and this chapter states it rather than hides it:
\textbf{the decision quantity lives in the arm whose uncertainty was just shown to be the least
trustworthy}. $\Prob(\tau > \text{cost})$ is read off the same posterior whose total-lift interval
covers at \pc{\nbSevenTotalCov}, and it inherits that optimism directly. The test is easy to run:
re-centre nothing, merely rescale the posterior's standard deviation to the design-based one, and
watch the probability move. It falls from \nbSevenPGo{} to \nbSevenPGoRescaled{}. In this run the
verdict survives --- \nbSevenDecision{} either way --- but that is luck, not robustness. A run
sitting near the decision bar would have had its rollout call overturned by the same correction,
and nothing in the posterior's own output would have warned anyone.

The ledger, then, stated plainly. Bayes did not buy the point estimate: the classical one is at
least as good. It did not buy interpretable weights: the classical blend is sparser. It did not
buy a trustworthy interval on the total: it bought a worse one. What it bought is real but
narrower than the brochure --- a near-nominal per-week band, a coherent way to propagate
uncertainty into euros, and the only object in the chapter that can answer the question the
business actually asked. \textbf{Use the posterior for the decision, and design-based inference
for the confidence in it.}

% =====================================================================================
\section{The decision, in euros}
\label{sec:sc:euro}

The campaign cost \keur{\nbSevenCost}. Everything above now collapses into one object: the
posterior over $\tau$. The rollout question is not ``is the lift positive?'' --- the placebo test
of \cref{sec:sc:classical} settled that --- but \textbf{``what is the probability the lift beats
what was paid?''}, because a real-but-small effect still loses money at national scale.

The decision rule is a stated convention, applied to the posterior:
%
\begin{equation}
  \Prob(\tau > c)\ \begin{cases}
    > 0.9 & \Rightarrow\ \textsc{go} \quad\text{(act without buying more evidence)},\\
    < 0.5 & \Rightarrow\ \textsc{no-go} \quad\text{(more likely than not it does not pay)},\\
    \text{otherwise} & \Rightarrow\ \textsc{test further}.
  \end{cases}
  \label{eq:sc:rule}
\end{equation}
%
At $c = \keur{\nbSevenCost}$ the posterior gives $\Prob(\tau > c) = \nbSevenPGo$, an expected
total of \keur{\nbSevenExpValue}, an expected net value of \keur{\nbSevenExpNet}, and an expected
return on investment of \pc{\nbSevenExpRoi}. The verdict is \textbf{\nbSevenDecision}
(\cref{fig:sc:decision}).

\input{build/floats/nb07_euro}

Two things are true at once here, and conflating them is how marginal campaigns get scaled. The
placebo test says the effect is \emph{real} ($p = \nbSevenPRatio$). The economics say it is
\emph{marginal at this price}: the expected lift sits almost exactly on the cost line, and
\nbSevenPGo{} is barely better than a coin flip. Statistical significance and commercial
worthiness are different questions asked of different objects.

The right-hand panel of \cref{fig:sc:decision} prices the gap. Sweeping a hypothetical cost $c$
and plotting $\Prob(\tau > c)$ gives a headroom curve, and its crossing of the $0.9$ bar sits at
\keur{\nbSevenHeadroom}. Had the campaign been bought for that instead of \keur{\nbSevenCost},
this same evidence would already be a confident \textsc{go}. That is a number to take into the
media-buying negotiation, and it is a more useful deliverable than a yes or a no.

\subsection*{What a confirmatory test is worth}

\textsc{test further} must not end as a shrug. The deliverable is an answer to \emph{what evidence
would settle this, and is it worth buying?} The value of perfect information,
%
\begin{equation}
  \text{VOI} \;=\; \E\big[\max(\tau - c,\ 0)\big] \;-\; \max\big(\E[\tau] - c,\ 0\big),
  \label{eq:sc:voi}
\end{equation}
%
is the expected gain from deciding \emph{after} the uncertainty resolves rather than committing
now on the mean \citep{raiffa1961}. It comes to \keur{\nbSevenVoi} --- and that small number is
itself the finding. With the posterior centred almost exactly on the cost line, even perfect
foresight barely changes the expected outcome, because whichever way the truth falls the firm is
near break-even. \Cref{eq:sc:voi} is therefore an upper bound on any confirmatory-research
budget: a study costing many times \keur{\nbSevenVoi} to resolve a \keur{\nbSevenVoi} question is
not a cautious decision, it is an expensive one.

Sizing that test sharpens the point. A second geo test treating $k$ independent markets tightens
the posterior roughly as $1/\sqrt{k}$; the smallest $k$ that reaches
$\Prob(\tau > c) \ge 0.9$ is \nbSevenKNeeded{} --- more treated markets than this world contains
(\nbSevenNDmas{} in total, one already used). And that arithmetic \emph{flatters} the test twice
over: it treats the $k$ market-level estimates as independent when they share the same macro
factors, and it takes the per-market scale from the posterior that \cref{sec:sc:compare} showed to
be too tight. Read \nbSevenKNeeded{} as a lower bound. The honest conclusion is that no feasible
single confirmatory geo test settles this question at this price. What remains is to renegotiate
the price (the headroom curve says \keur{\nbSevenHeadroom} clears the bar), to extend the
post-period, or to accept a lower confidence bar with open eyes. What is not on the list is a
blind national rollout.

% =====================================================================================
\section{What can go wrong}
\label{sec:sc:wrong}

\subsection*{The naive estimators}

\Cref{tab:sc:naive} runs the shortcuts of \cref{sec:sc:decision} on the same panel. Before-and-after
books the trend and overstates the weekly lift badly. Treated-minus-average-control is
\emph{worse}, because it differences nothing out and simply books the level gap between a large
metro and the average one. Differencing out the shared pre-period level --- difference-in-differences
against the average control --- removes the trend bias, but a residual gap survives, because the
average control has the wrong factor loadings and so does not cancel the common wave cleanly. Only
the reweighted control lands near the truth. Same data, four estimates: the estimator \emph{is} the
identification strategy.

\input{build/tables/nb07_naive}

\subsection*{Anticipation}

\Cref{as:sc:anticip} is tested by moving the launch backwards. Refitting with a fake launch in week
\nbSevenPlaceboTimeWeek{} --- inside the true pre-period, where by construction nothing happened ---
must return an effect of approximately zero, and does: \keur{\nbSevenPlaceboTimeGap} per week,
comfortably inside the pre-fit noise of \keur{\nbSevenBayesPreRmse}. The method does not invent
lift where there is none.

\subsection*{Spillover, the assumption no placebo can see}

\Cref{as:sc:sutva} is the one that placebo inference structurally cannot test: if the campaign
leaks into neighbouring markets, then every placebo refit is contaminated in exactly the same way
and nothing looks unusual. The simulator can do what no real dataset can, and plant the violation.
A fraction $\varphi$ of the true lift is leaked into \nbSevenSpNNeighbours{} neighbouring donors,
$Y_{jt} \mathrel{+}= \varphi\,\Delta_t$, and the estimator is refit.

The first-order prediction, before running it, is that the contaminated donors rise, so the
synthetic rises, so the measured gap \emph{shrinks} by roughly the spilled weight:
%
\begin{equation}
  \text{bias} \;\approx\; -\,\varphi \sum_{j \in \text{neighbours}} w_j \;\times\; \tau .
  \label{eq:sc:spill}
\end{equation}
%
At $\varphi = \nbSevenSpPhi$, with \pc{\nbSevenSpNeighbourWeight} of the weight on the
contaminated donors, \cref{eq:sc:spill} predicts a bias of $\nbSevenSpPredicted$ (EUR 000) and
the refit delivers $\nbSevenSpMeasured$ (\cref{tab:sc:spillover}).

\input{build/tables/nb07_spillover}

Two consequences follow, and the first is more reassuring than most analysts expect.
\textbf{Spillover attenuates}: synthetic control \emph{under}-states a real campaign, never
inflates it. It is conservative for a \textsc{go} call --- though it can certainly turn a true
\textsc{go} into a \textsc{no-go}. And because no placebo can flag it, \textbf{the fix is design,
not statistics}: excluding the suspect donors up front restores the estimate to
\keur{\nbSevenSpExcluded}, against \keur{\nbSevenSpContaminated} for the contaminated pool and
\keur{\nbSevenSpClean} for the clean one. A smaller clean donor pool beats a larger contaminated
one. On real data, drop adjacent and media-overlapping markets from the pool before fitting
anything.

\subsection*{The limits}

Three remain, and they are properties of the design rather than defects of the code. \emph{One
treated unit buys limited power}: with \nbSevenNDmas{} markets, the permutation $p$-value is
granular and cannot resolve below roughly $1/\nbSevenNDmas$ however strong the effect. \emph{The
pre-fit gate is a gate}: a high pre-RMSE makes the post gap uninterpretable, and no amount of
Bayesian machinery repairs a bad donor pool. And \emph{the estimator is high-variance here}, as
\cref{sec:sc:valid} measured --- which is the reason this chapter leads its conclusions with the
placebo test rather than with either interval.

% =====================================================================================
\section{On real data}
\label{sec:sc:real}

Everything above was graded against a truth that was planted. The obvious objection --- that a
method which recovers a lift someone put there is not thereby proved to work --- deserves a real
answer, and there is one. A real advertiser spent \kusd{\nbSevenBSpend} on a real campaign in
\nbSevenBNTreated{} US geos over \nbSevenBCampaignDays{} days in early 2015, leaving
\nbSevenBNControl{} geos dark; the daily sales-and-cost panel is published with Google's geo
experiment methodology \citep{vaver2011, kerman2017}. It is exactly the shape of panel every
advertiser's warehouse already holds: \nbSevenBNGeos{} units, \nbSevenBNDaysTotal{} days,
one campaign window, and \nbSevenBCooldownDays{} days of cooldown after the money stops.

Three things replace the planted truth.

\subsection*{The referee: what randomization buys}

The \nbSevenBNGeos{} geos were split into treatment and control \emph{at random}, in
\nbSevenBNPairs{} size-matched pairs. That design licenses an estimator that fits no model at
all. Take each pair's treated geo, take its change in mean daily sales from the pre-window to the
campaign window, subtract the same change for its matched control partner, and average over the
pairs:
%
\begin{equation}
  d_p \;=\; \big(\bar Y^{\,\text{tr}}_{p,\text{post}} - \bar Y^{\,\text{tr}}_{p,\text{pre}}\big)
          - \big(\bar Y^{\,\text{ct}}_{p,\text{post}} - \bar Y^{\,\text{ct}}_{p,\text{pre}}\big),
  \qquad
  \hat\tau_{\text{DiD}} \;=\; N T \cdot \frac{1}{P}\sum_{p=1}^{P} d_p ,
  \label{eq:sc:did}
\end{equation}
%
with $N = \nbSevenBNTreated$ geos, $T = \nbSevenBCampaignDays$ days and $P = \nbSevenBNPairs$
pairs, and a standard error built from the scatter of the pair differences,
$NT \cdot s_d / \sqrt{P}$. The covariance choice is not a default: randomization happened at the
geo level within pairs, so the pair differences are the design's own unit of variability. A
geo-day-level iid standard error would treat thousands of correlated geo-days as independent
draws and report an interval several times too narrow --- \citeauthor{bertrand2004}'s warning,
one chapter later and in a different costume. Pairing is itself variance reduction: the geo size
factor (the largest geo sells \nbSevenBSizeRatio{} times the smallest) differences out
\emph{within} each pair, and ignoring the pairs inflates the standard error from
\kusd{\nbSevenBDidSe} to \kusd{\nbSevenBUnpairedSe}, a factor of \nbSevenBPairingGain{}.

\Cref{eq:sc:did} is \emph{unbiased by construction}. It is the closest thing to truth that real
data ever hands anyone, and it is what replaces the planted lift: every model-based estimate must
land inside its interval. It gives \kusd{\nbSevenBDidTotal} of incremental sales, 90\,\% interval
$[\kusd{\nbSevenBDidLo},\ \kusd{\nbSevenBDidHi}]$, an incremental return on ad spend of
\nbSevenBDidIroas{}. Note the width. Unbiasedness is not precision, and that width --- nearly
half the effect in each direction --- is the entire commercial reason to fit a model on top of a
perfectly good experiment.

\subsection*{The contestant, and the exam}

The classical synthetic control of \cref{sec:sc:classical}, applied unchanged --- the treated unit
being the \emph{cell}, the mean daily sales per treated geo, so that the target stays inside the
donors' convex hull (\cref{as:sc:hull}) --- tracks the cell to a pre-fit RMSE of
\usd{\nbSevenBScPreRmse} per geo per day on a level of \usd{\nbSevenBPreLevel}, or
\pc{\nbSevenBScPreRmsePct} (\cref{fig:sc:realsc}). It returns \kusd{\nbSevenBScTotal} ---
$\nbSevenBScErrVsAnchor$ (USD 000) from the randomized referee, or \nbSevenBScAnchorSes{} anchor
standard errors, comfortably inside its interval. The Bayesian fit of \cref{sec:sc:bayes} returns
\kusd{\nbSevenBBayesTotal}, \nbSevenBBayesAnchorSes{} anchor standard errors from the referee.
\Cref{tab:sc:realcompare} records all three.

\input{build/floats/nb07b_sc}
\input{build/tables/nb07b_compare}

That is the exam, and it is the only exam real data can set: the method passes it. And the same
bake-off runs, with no simulator, on a single real market --- pretending only one mid-size geo had
been treated reproduces \cref{tab:sc:naive}'s failures exactly (\cref{tab:sc:realnaive}), including
one the simulation could only warn about. The largest geo in the country is bigger than every
control geo and therefore lies \emph{outside} its own donor hull. \Cref{as:sc:hull} fails, the
simplex does its honest best by piling all its weight on the single largest donor, and the pre-fit
gate \textbf{fails loudly}: a pre-RMSE of \usd{\nbSevenBHullPreRmse} against a bar of
\usd{\nbSevenBHullGateRef}. The estimate it \emph{would} have shipped implies an iROAS of
\nbSevenBHullIroas{} against the experiment's \nbSevenBDidIroas{} --- precisely the inflated
number an unscrupulous deck would have quoted. The method's own diagnostic refuses it. That
self-diagnosis is worth more than the estimate.

\input{build/tables/nb07b_naive}

\subsection*{The off-switch}

Real data affords one falsification that no simulator provides for free: \textbf{the campaign
ends}. The spend stops, and a genuine advertising effect must collapse when the money does. It
does (\cref{fig:sc:realgap}): a sustained $\nbSevenBPostGap$ dollars per geo per day through the
paid window, and $\nbSevenBCoolGap$ dollars per geo per day in the \nbSevenBCooldownDays-day
cooldown. An effect that persisted after the spend stopped would smell of a confound rather than a
campaign. This one switches off --- which simultaneously certifies the attribution and tells the
media planner there is no carryover annuity to book. The placebo apparatus of
\cref{sec:sc:classical} agrees: on \nbSevenBNPseudo{} pseudo-cells of untreated geos, the real
cell is the most extreme unit observed ($p = \nbSevenBPRatio$ on the RMSE ratio,
$p = \nbSevenBPGap$ on the raw gap), and a fake launch inside the pre-period finds
$\nbSevenBPlaceboTimeGap$ dollars per geo per day --- \pc{\nbSevenBPlaceboTimePct} of the real
effect, and the wrong sign.

\input{build/floats/nb07b_gap}

\subsection*{The interval, re-priced --- the chapter's finding, confirmed without a simulator}

Fresh worlds are not available here, so coverage cannot be \emph{measured}. But it need not be
guessed either. The placebos re-run \emph{the same estimator} under a known null, and that is the
fair referee for the posterior's width. \Cref{tab:sc:realscales} sets the scales against each
other. The posterior's standard deviation is \kusd{\nbSevenBBayesSd}. The placebo distribution of
the same estimator gives \kusd{\nbSevenBPlaceboSd}, and even after deflating it for the fact that
a half-size pseudo-cell is noisier than the real cell --- the conservative floor ---
\kusd{\nbSevenBPlaceboSdFloor}. The posterior is too tight by
\nbSevenBRatioPlaceboFloor--\nbSevenBRatioPlacebo$\times$.

\input{build/tables/nb07b_scales}

That is \cref{sec:sc:compare}'s finding, reproduced on real data, in the same direction, with the
same cause. What was an imported claim from a simulator is now a measured one.

One thing this does \emph{not} say, and the distinction is worth the sentence. The randomized
anchor's standard error is \nbSevenBRatioAnchor$\times$ the posterior's, and that ratio is
\textbf{not} the indictment. An unweighted matched-pairs difference-in-differences is a genuinely
noisier estimator than a donor-matched synthetic control; cutting that variance is the entire job
of the model, and it really does cut it. The fair referee for the posterior's width is the placebo
distribution of the \emph{same} estimator, and that referee returns a verdict of roughly two, not
eight. Both statements hold at once: \textbf{synthetic control is far more precise than the raw
experiment, and the posterior is still overconfident about how precise.}

\subsection*{The dollar decision, and the margin trap}

The trap in every iROAS deck: an iROAS of three is not a threefold return. Ad spend buys
\emph{revenue}; the firm keeps only the gross margin on it. The campaign pays if
%
\begin{equation}
  \text{iROAS} \times \text{margin} \;>\; 1
  \qquad\Longleftrightarrow\qquad
  \text{iROAS} \;>\; \frac{1}{\text{margin}} ,
  \label{eq:sc:margin}
\end{equation}
%
so at the assumed \pc{\nbSevenBMargin} gross margin the bar is not $1.0$ but
\nbSevenBBreakEvenIroas{} --- and the posterior sits almost exactly on top of it. The break-even
margin is $1/\text{iROAS} = \pc{\nbSevenBBreakEvenMargin}$, and the campaign reaches 90\,\%
confidence of profitability only above \pc{\nbSevenBMarginNinety} (\cref{fig:sc:realmargin}). The
lift is beyond doubt; the profitability turns entirely on a finance input the model cannot see.

\input{build/floats/nb07b_margincurve}

And the tension of \cref{sec:sc:compare} arrives here too, in dollars. At the posterior's own
scale, $\Prob(\text{profitable}) = \nbSevenBPProfit$. Re-centre nothing and merely re-price at the
placebo floor and it falls to \nbSevenBPProfitFloor{}; at the raw placebo scale,
\nbSevenBPProfitRaw{}; at the anchor's, \nbSevenBPProfitAnchor{} (\cref{tab:sc:realscales}, last
column). The call survives at every scale --- \nbSevenBDecision{}, because at a
\pc{\nbSevenBMargin} margin the campaign is on a knife-edge that no re-scaling rescues --- but the
confidence attached to it was borrowed, not earned.

\begin{remark}[The lecture point]
The randomized experiment did not merely validate a number. It validated \emph{the method that
will be used when there is no experiment}. The moment a firm has one treated market chosen by a
media planner rather than a coin, everything in this section except the anchor disappears --- and
synthetic control is the only member of the family still standing. That is why the chapter grades
it here, where the answer is known.
\end{remark}

% =====================================================================================
\section{Takeaways}
\label{sec:sc:take}

\begin{enumerate}[leftmargin=*, itemsep=.45em]
  \item \textbf{The counterfactual is built, not found.} With no control group, the estimator
        \emph{is} the identification strategy. A weighted blend of donors, constrained to the
        simplex so it cannot extrapolate, reconstructs the treated market's untreated path from
        the factors it shares with everyone else. The naive shortcuts do not fail by a little
        (\cref{tab:sc:naive}); they fail by booking the shared macro wave as campaign effect.

  \item \textbf{The pre-fit is a gate, not a formality.} If the synthetic cannot track the
        treated unit \emph{before} launch, the post-launch gap means nothing, and no Bayesian
        machinery repairs a bad donor pool. On real data the gate is what stops the country's
        largest market --- outside its own donor hull --- from shipping an iROAS of
        \nbSevenBHullIroas{} against a true \nbSevenBDidIroas{}.

  \item \textbf{Classical and Bayesian agree on the number.} They differ by
        \pc{\nbSevenEstDiffPct} of the effect, and the classical point estimate is, if anything,
        the better of the two. The sampler adds no causal content; the identification does all of
        it.

  \item \textbf{The Bayesian interval on the total under-covers badly --- and the reason is a
        misspecified likelihood, not the paradigm.} A nominal 90\,\% credible interval covers on
        \pc{\nbSevenTotalCov} of fresh panels, against \pc{\nbSevenRiCov} for the design-based
        randomization interval. The posterior prices the standard deviation of the
        \nbSevenNPost-week total at \keur{\nbSevenBayesSd} when the truth is
        \keur{\nbSevenSeedSd}. \Cref{eq:sc:n2} is the whole story: the errors are a persistent
        offset, the variance of a sum under near-perfect autocorrelation approaches $n^2\sigma^2$
        rather than $n\sigma^2$, and \cref{eq:sc:lik} set every autocorrelation to zero. Repair
        the likelihood (correlated errors) or repair the posterior
        \citep[sandwich-corrected;][]{mueller2013}. Do not repair your faith in Bayes.

  \item \textbf{The classical robustness toolkit wins here because it committed to less.}
        Randomization inference never wrote a likelihood, so it had no likelihood to get wrong; it
        measures the estimator's error distribution from units known to be untreated, persistent
        macro mismatch and all. That is not a defeat for Bayesian inference. It is a reminder of
        what a design-based procedure is \emph{for}.

  \item \textbf{Only the posterior can answer the question the business asked.}
        $\Prob(\tau > \text{cost})$ has no classical counterpart, and \cref{eq:sc:rule} consumes
        nothing else. Hence the discipline this chapter recommends, and the sentence the book will
        keep returning to: \textbf{use the posterior for the decision, and design-based inference
        for the confidence in it} --- and remember that the euro probability inherits whatever
        optimism the interval carries.

  \item \textbf{Significant and profitable are different questions.} The effect here is real
        ($p = \nbSevenPRatio$) and marginal at the price paid ($\Prob(\tau > c) = \nbSevenPGo$,
        expected net value \keur{\nbSevenExpNet}). Conflating the two is how marginal campaigns
        get scaled nationally. The headroom curve --- this evidence clears the bar at a price of
        \keur{\nbSevenHeadroom} --- is the more useful deliverable than either a $p$-value or a
        verdict.
\end{enumerate}

\notebooknote{%
  The analysis of \crefrange{sec:sc:dgm}{sec:sc:wrong} is
  \code{notebooks/07\_geo\_lift\_synthetic\_control.ipynb};
  \cref{sec:sc:real} is \code{notebooks/07b\_geo\_lift\_real\_data.ipynb}, which loads Google's
  published geo-experiment panel \citep{kerman2017} directly. Both run in a fast teaching mode
  (\code{CMP\_FAST=1}, short chains) and a full mode (\code{CMP\_FAST=0}); every number in this
  chapter comes from the full run, emitted by \code{cmp.report} and injected here as a macro. The
  simulator is \code{cmp.dgp.geo\_panel}, the estimators are \code{cmp.estimators}, and the seeds
  are fixed in the notebooks. \citeauthor{abadie2010}'s California Proposition~99 study remains the
  canonical public panel on which to practise the method \citep{abadie2010}.}
